% Ajustes de composición de la memoria técnica de HispaLIS.
% Se pasa a pandoc con --include-in-header desde docs/generar-memoria-pdf.sh.

\usepackage{etoolbox}

% Las tablas de este documento son anchas (la de R5 frente a R4, la de los 44 ADR).
% Un cuerpo menor dentro de longtable es lo que evita que se salgan del margen.
\AtBeginEnvironment{longtable}{\small}

% Los nombres de fichero de los ADR son largos y van en tipografía de máquina de
% escribir, que LaTeX no parte por sílabas: `\ttfamily` deja el carácter de guion
% desactivado y una palabra más ancha que la columna se sale del margen sin remedio.
% Esto lo vuelve a activar en cuanto se selecciona la tipografía. Es lo que hace el
% paquete `hyphenat` con su opción `htt`, escrito a mano porque la distribución de TeX
% de la imagen de pandoc no lo trae — y depender de un paquete que puede no estar es
% justo lo que un guion de publicación no debe hacer.
\makeatletter
\g@addto@macro\ttfamily{\hyphenchar\font=`\-\relax}
\makeatother

% Los guiones explícitos también son puntos de corte válidos, y aquí interesan.
\exhyphenpenalty=10
\hyphenpenalty=50

% Margen de maniobra para el ajuste de línea en párrafos con muchas URLs y código.
% Antes una línea demasiado holgada que una que invade el margen.
\tolerance=2000
\setlength{\emergencystretch}{4em}

% Las figuras van donde se las pone: son parte del hilo del texto, no ilustraciones
% que puedan flotar tres páginas más allá de lo que explican.
\usepackage{float}
\floatplacement{figure}{H}
